\documentclass[sigconf,nonacm]{acmart}

% ---------------------------------------------------------------
% Packages
% ---------------------------------------------------------------
\usepackage{listings}
\usepackage{xcolor}
\usepackage{booktabs}
\usepackage{graphicx}
\usepackage{hyperref}
\usepackage{microtype}
\usepackage{amsmath}
\usepackage{balance}

% ---------------------------------------------------------------
% Listings style
% ---------------------------------------------------------------
\definecolor{codebg}{HTML}{F7F7F7}
\definecolor{codekey}{HTML}{0000AA}
\definecolor{codestr}{HTML}{AA0000}
\definecolor{codecmt}{HTML}{007700}

\lstdefinelanguage{GraphQL}{
  keywords={query, mutation, type, input, enum, interface, scalar,
            on, fragment, subscription, true, false, null},
  keywordstyle=\color{codekey}\bfseries,
  stringstyle=\color{codestr},
  commentstyle=\color{codecmt}\itshape,
  morestring=[b]",
  morestring=[b]',
  morecomment=[l]{\#},
  sensitive=true,
}

\lstset{
  backgroundcolor=\color{codebg},
  basicstyle=\small\ttfamily,
  breaklines=true,
  frame=single,
  framesep=4pt,
  rulecolor=\color{gray!40},
  captionpos=b,
  numbers=left,
  numberstyle=\tiny\color{gray},
  numbersep=6pt,
  xleftmargin=14pt,
}

% ---------------------------------------------------------------
% Metadata
% ---------------------------------------------------------------
\title{Mozg: A ``Slow by Design'' Federated Query Layer\\
       for Heterogeneous Data Sources via GraphQL}

\author{Dominik Maszczyk}
\affiliation{%
  \institution{Independent Research}
  \country{}}
\email{skitionek@gmail.com}

\begin{abstract}
Modern data ecosystems fragment knowledge across relational databases,
graph stores, document databases, REST services, and domain-specific
repositories.  Integrating these silos typically requires either a costly
centralised data warehouse or bespoke per-integration middleware.
We present \textbf{Mozg} (Polish: \emph{brain}), an open-source
federated query layer that exposes a single GraphQL endpoint capable
of routing queries to any supported backend at runtime, traversing
user-defined cross-source relationships, and synthesising results
without ever storing the underlying data.  Mozg deliberately trades
query latency for data freshness, flexibility, and minimal operational
overhead---a philosophy we call \emph{``slow by design''}.
We describe the architecture, driver model, ontology-ingestion pipeline,
and the browser-based visual query builder.  We evaluate the system
against representative public databases and REST APIs, discuss
scalability and security considerations, and position Mozg relative
to related federated-query systems.
\end{abstract}

\keywords{federated queries, GraphQL, heterogeneous data integration,
          ontology mapping, data virtualisation, cross-database joins}

% ---------------------------------------------------------------
\begin{document}
\maketitle
% ---------------------------------------------------------------

% ===============================================================
\section{Introduction}
\label{sec:intro}
% ===============================================================

The proliferation of specialised data stores---relational databases,
property-graph engines, triple stores, document databases, column
stores, and REST services---creates what is commonly called the
\emph{data silo} problem~\cite{halevy2006data}.  Organisations that
need to ask questions spanning multiple silos face an uncomfortable
choice: build a data warehouse that periodically copies and integrates
data, or write custom integration code for each pair of sources.
Both options introduce significant cost: warehouses require ETL
pipelines, transformation logic, and storage duplication; bespoke code
is hard to maintain and becomes stale the moment source schemas evolve.

\textbf{Mozg} takes a third approach.  Instead of centralising data, it
centralises the \emph{ability to query} data.  Users describe the
connections between entities in different systems---a
\texttt{hasMany}/\texttt{hasOne}/\texttt{belongsTo} relationship lattice
that mirrors familiar ORM conventions---and Mozg resolves these
relationships at query time, returning a joined result without ever
persisting a copy of the underlying records.  The single entry-point
is a GraphQL endpoint; users supply their own database credentials
as part of each query, keeping the server stateless with respect to
authorisation.

The name alludes to the system's role as a \emph{cognitive layer}: like
a brain, Mozg does not store memories---it connects them.  The
``slow by design'' label is not an apology but a deliberate design
decision: latency is accepted as the price of freshness and schema
flexibility.  This trade-off is appropriate for interactive exploration,
scientific data synthesis, and ad-hoc cross-domain analytics where
staleness is unacceptable.

The remainder of this paper is organised as follows.
Section~\ref{sec:background} reviews related work.
Section~\ref{sec:architecture} presents the overall architecture.
Section~\ref{sec:drivers} describes the driver model.
Section~\ref{sec:ontology} covers the OWL ontology pipeline.
Section~\ref{sec:ui} describes the query-builder interface.
Section~\ref{sec:evaluation} reports an empirical evaluation.
Section~\ref{sec:discussion} discusses limitations and future work.
Section~\ref{sec:conclusion} concludes.

% ===============================================================
\section{Background and Related Work}
\label{sec:background}
% ===============================================================

\subsection{Federated Query Systems}

Federated databases were formalised by Sheth and Larson~\cite{sheth1990federated},
who classified multi-database systems along axes of autonomy and
heterogeneity.  Classical systems such as IBM's Garlic~\cite{haas1997optimizing}
and TSIMMIS~\cite{hammer1995tsimmis} translate a unified query language
into source-specific sub-queries.  More recent systems include Presto
(now Trino)~\cite{sethi2019presto}, which federates analytical SQL over
HDFS, Hive, and relational stores, and Apache Drill~\cite{hausenblas2013apache},
which enables schema-free SQL across Hadoop and NoSQL backends.

\subsection{GraphQL as a Query Interface}

GraphQL~\cite{graphql2015} was introduced by Facebook in 2015 as a
typed query language for APIs.  Its hierarchical, declarative nature
maps well to the entity--relationship model.  GraphQL Federation
(Apollo Federation)~\cite{apollofederation} distributes a single
GraphQL schema across multiple microservices and sub-graphs.  Unlike
Apollo Federation, which federates GraphQL \emph{services},
Mozg federates \emph{raw data sources} that are not themselves GraphQL
services.

GraphQL Mesh~\cite{graphqlmesh} provides source handlers that translate
a unified GraphQL schema into REST, gRPC, SOAP, OpenAPI and database
calls.  Mozg uses GraphQL Mesh handlers for most relational and API
backends while retaining custom drivers for specialised sources
(ArangoDB, BioCyc) and providing its own relation-traversal layer on
top.

\subsection{Data Virtualisation}

Data virtualisation~\cite{datavirt} provides a logical, unified view of
distributed data without physical integration.  Commercial products such
as Denodo, TIBCO Data Virtualization, and IBM Federation Server occupy
this space.  Mozg is a lightweight, open-source alternative positioned
for exploratory use: it does not implement query optimisation, query
pushdown strategies, or persistent schema registries.

\subsection{Semantic Web and Ontology Mapping}

OWL (Web Ontology Language)~\cite{owl2} and RDFS define vocabularies and
axioms for machine-readable knowledge.  Systems such as
Ontop~\cite{ontop} and D2RQ~\cite{d2rq} map OWL/SPARQL queries onto
relational databases (SPARQL-to-SQL rewriting).  Mozg's ontology
pipeline moves in the opposite direction: given an OWL TBox, it
generates GraphQL type definitions and entity maps that can then be
used to drive its federated query layer.

% ===============================================================
\section{Architecture}
\label{sec:architecture}
% ===============================================================

\begin{figure}[h]
\centering
\includegraphics[width=0.95\columnwidth]{mozg-architecture}
\caption{High-level architecture of Mozg.}
\label{fig:arch}
\end{figure}

Figure~\ref{fig:arch} shows the five-layer architecture of Mozg.

\paragraph{HTTP Server (\texttt{src/index.js}).}
A Node.js server built on \emph{graphql-yoga}~\cite{graphqlyoga}
listens on a configurable port (default 4000).  It exposes three
URL namespaces: \texttt{/graphql} for the GraphQL endpoint and
GraphiQL IDE, \texttt{/} for the browser query builder, and
\texttt{/examples} for downloadable ontology samples.  A lightweight
static-file handler serves the \texttt{public/} directory without
external middleware.

\paragraph{GraphQL Schema (\texttt{src/schema.js}).}
A hand-written SDL schema defines the full type system.  There are
three root fields:
\begin{itemize}
  \item \texttt{query(input: QueryInput!)} --- executes a federated
        multi-source query with optional relation traversal;
  \item \texttt{introspect(connection)} --- returns the schema of a
        remote database (tables, columns, types, nullability,
        primary-key flags);
  \item \texttt{catalog(name)} --- returns one or all entries from
        the curated database catalog.
\end{itemize}
A single mutation, \texttt{ingestOntology(input)}, parses an OWL
ontology and returns extracted classes, properties, and generated
GraphQL SDL.  A custom \texttt{JSON} scalar accepts arbitrary
structured values so that heterogeneous query results can be
returned without a fixed schema.

\paragraph{Query Connector (\texttt{src/database/connector.js}).}
A thin routing layer resolves the \texttt{driver} field of the
supplied \texttt{ConnectionInput} via the driver registry and
delegates execution.

\paragraph{Driver Registry (\texttt{src/database/registry.js}).}
The single source of truth for driver bindings.  It maps the thirteen
driver names (\texttt{postgres}, \texttt{mysql}, \texttt{sqlite3},
\texttt{neo4j}, \texttt{arango}, \texttt{biocyc}, \texttt{rest},
\texttt{kegg}, \texttt{openapi}, \texttt{soap}, \texttt{odata},
\texttt{thrift}, \texttt{mongodb}) to their implementing modules.

\paragraph{Driver Implementations (\texttt{src/database/drivers/}).}
Each driver exposes a uniform interface: \texttt{executeQuery(input)}
and, optionally, \texttt{introspect(connection)}.  Most drivers are
thin wrappers around \emph{GraphQL Mesh} source handlers; three
legacy/specialised drivers (SQLite3 via Knex, ArangoDB, BioCyc) use
custom client libraries.

\subsection{Query Execution Flow}

When a client issues a \texttt{query} request, the server:
\begin{enumerate}
  \item Parses the \texttt{QueryInput} (connection, table/entity,
        columns, predicates, limit/offset/orderBy, relations).
  \item Resolves the driver and calls \texttt{executeQuery}.
  \item The driver issues the primary query against the target source.
  \item For each \texttt{RelationInput} in the relations array, the
        driver issues one sub-query per relation, keyed on the
        foreign-key values returned by the primary query.
  \item Results are merged and returned as a \texttt{QueryResult}
        with a \texttt{count} and a JSON \texttt{data} array.
\end{enumerate}

This is a straightforward nested-loop join executed in application
code.  The design deliberately avoids query pushdown and cross-source
optimisation; it is simple to implement, easy to debug, and
independent of any particular source's query planner.  Relation
definitions can be nested to arbitrary depth, enabling multi-hop
traversals across heterogeneous sources.

\subsection{Connection Model}

Mozg is deliberately \emph{stateless} with respect to credentials.
All connection parameters (host, port, database name, username,
password, authentication scheme, HTTP headers) are supplied by the
client on every request via \texttt{ConnectionInput}.  The server
maintains an in-memory connection cache keyed on host and username
to avoid re-establishing connections within a session, but it does
not persist credentials to disk or a secrets store.  This design
simplifies the threat model: the server is not a credential vault and
does not require secret-management infrastructure.

% ===============================================================
\section{Driver Model}
\label{sec:drivers}
% ===============================================================

\subsection{GraphQL Mesh Adapter}

The \texttt{mesh-adapter.js} driver provides a unified implementation
for the eight source types handled by GraphQL Mesh handlers:
\texttt{postgres} (\texttt{@graphql-mesh/postgraphile}),
\texttt{mysql} (\texttt{@graphql-mesh/mysql}),
\texttt{neo4j} (\texttt{@graphql-mesh/neo4j}),
\texttt{openapi} (\texttt{@graphql-mesh/openapi}),
\texttt{soap} (\texttt{@graphql-mesh/soap}),
\texttt{odata} (\texttt{@graphql-mesh/odata}),
\texttt{thrift} (\texttt{@graphql-mesh/thrift}), and
\texttt{mongodb} (\texttt{@graphql-mesh/mongoose}).

Each Mesh handler generates its own schema from the remote source's
metadata (e.g., PostgreSQL information schema, OpenAPI spec, WSDL).
The adapter wraps this schema to present the uniform
\texttt{executeQuery}/\texttt{introspect} interface expected by the
connector layer.  Handler instances are cached per unique connection
string to avoid repeated metadata fetches.

\subsection{SQLite3 Driver}

The legacy SQLite3 driver uses \emph{Knex}~\cite{knex} for query
building and connection management.  All user-supplied identifiers
(table names, column names) are passed as Knex identifiers, and all
predicate values are bound as parameterised query arguments.  This
ensures that SQL injection via value parameters is not possible
(identifier injection via maliciously crafted table/column names
remains a known limitation tracked as a deferred issue).

\subsection{ArangoDB Driver}

The ArangoDB driver uses the official \emph{arangojs}~\cite{arangojs}
client.  It translates \texttt{QueryInput} fields into AQL (ArangoDB
Query Language) \texttt{FOR}/\texttt{FILTER}/\texttt{SORT}/\texttt{LIMIT}
clauses, binding all predicate values through AQL bind parameters.

\subsection{REST Driver}

The plain REST driver maps \texttt{QueryInput} fields to HTTP GET
requests against a base URL.  The entity name is appended as a path
segment; \texttt{where} predicates are converted to query-string
parameters.  Because there is no schema metadata, the driver returns
the raw JSON response.  A known limitation is the N+1 sub-request
pattern for relation traversal; each related entity produces one HTTP
request per primary result row rather than a batched call.  This is
tracked as a deferred performance issue.

\subsection{BioCyc Driver}

The BioCyc driver integrates with the BioCyc~\cite{biocyc} collection
of metabolic pathway databases via the BioCyc Web Services API.
It supports pathway, reaction, compound, and organism queries.  Because
the BioCyc API does not expose a machine-readable schema, the driver
relies on hard-coded response field mappings.

\subsection{KEGG Driver}

KEGG~\cite{kegg} (Kyoto Encyclopedia of Genes and Genomes) exposes its
REST API (\texttt{rest.kegg.jp}) as \texttt{text/plain} responses rather
than JSON.  The dedicated \texttt{kegg} driver transparently converts
these responses before they reach the query layer:

\begin{itemize}
  \item \textbf{Tab-delimited format} (used by \texttt{/list},
        \texttt{/find}, and \texttt{/link}) is parsed into row objects
        with fields \texttt{entry\_id}/\texttt{name} or
        \texttt{source\_id}/\texttt{target\_id} depending on the
        endpoint family.
  \item \textbf{KEGG flat-file format} (used by \texttt{/get}) is parsed
        into a single record object whose keys correspond to KEGG
        field names (e.g.\ \texttt{entry}, \texttt{formula},
        \texttt{pathway}) with repeated fields accumulated into arrays.
\end{itemize}

Because KEGG encodes query terms as URL \emph{path} segments rather
than query-string parameters (e.g.\ \texttt{/find/compound/glucose}),
the driver introduces a special \texttt{where.\_pathSuffix} convention:
any value supplied under that key is appended to the entity path
before the HTTP request is issued.  This allows callers to use
\texttt{/find/compound} as the entity name and pass
\texttt{where: \{\ \_pathSuffix: "glucose"\ \}} without hard-coding
the search term in the catalog.

% ===============================================================
\section{OWL Ontology Pipeline}
\label{sec:ontology}
% ===============================================================

A key distinguishing feature of Mozg is its ability to consume
OWL ontologies and convert them into queryable GraphQL schemas.
The pipeline has three stages: \emph{parsing}, \emph{extraction},
and \emph{mapping}.

\subsection{Parsing}

Mozg supports four serialisation formats:

\begin{itemize}
  \item \textbf{Turtle} (.ttl) --- parsed using the \emph{N3.js}
        library~\cite{n3js} via a streaming quad parser;
  \item \textbf{RDF/XML} (.rdf, .owl) --- parsed using
        \emph{rdfxml-streaming-parser}~\cite{rdfxmlparser};
  \item \textbf{OWL/XML} --- parsed using \emph{fast-xml-parser}~\cite{fxp}
        with a custom OWL/XML extractor;
  \item \textbf{Manchester Syntax} (.omn) --- parsed with a hand-written
        recursive-descent parser.
\end{itemize}

Format auto-detection inspects the HTTP \texttt{Content-Type} header,
file extension, and content heuristics (BOM, XML declaration,
Turtle prefix declarations) in that order.  Parsed quads are
normalised into a common internal representation before being passed
to the extraction stage.

\subsection{Extraction}

The extractor walks the quad set and collects:
\begin{itemize}
  \item \textbf{OWL Classes} --- IRI, rdfs:label, rdfs:subClassOf chains,
        and owl:deprecated flag;
  \item \textbf{Object Properties} --- IRI, label, domain, range, and
        an inferred \texttt{relationType}
        (\texttt{hasMany}/\texttt{hasOne}/\texttt{belongsTo}) based on
        cardinality restrictions;
  \item \textbf{Data Properties} --- IRI, label, domain, range (XSD type
        mapped to GraphQL scalar).
\end{itemize}

\subsection{Mapping}

The mapper converts the extracted TBox into a GraphQL SDL string and
an \emph{entity map}: a mapping from type names to metadata including
IRI, inferred table name, abstract flag, field list, and relation list.
Abstract OWL classes (those that have subclasses but no data properties
of their own) are emitted as GraphQL \texttt{interface} types;
concrete classes become \texttt{type} definitions.  XSD datatypes are
mapped to GraphQL scalars: \texttt{xsd:integer}/\texttt{xsd:int}
$\to$ \texttt{Int}; \texttt{xsd:decimal}/\texttt{xsd:double}/\texttt{xsd:float}
$\to$ \texttt{Float}; \texttt{xsd:boolean} $\to$ \texttt{Boolean};
all others $\to$ \texttt{String}.

\subsection{Caching}

To avoid re-parsing the same ontology repeatedly within a session,
the pipeline maintains an in-memory cache keyed on the SHA-256 hash
of the raw ontology content.  Cache hits return the stored result
immediately; misses trigger a full parse/extract/map cycle.

\subsection{Validation}

When the caller supplies a database \texttt{connection} alongside the
ontology input, Mozg optionally validates the generated schema against
the live database: it checks that every mapped class resolves to an
existing table and that every data property resolves to an existing
column.  Validation errors and warnings are returned in a structured
\texttt{ValidationReport}.

% ===============================================================
\section{Query Builder UI}
\label{sec:ui}
% ===============================================================

To lower the barrier to entry for non-technical users, Mozg ships a
browser-based visual query builder served at the root path
(\texttt{/}).  The interface allows users to:
\begin{enumerate}
  \item Select a driver and fill in connection parameters via a form;
  \item Pick a table or API endpoint from an auto-populated dropdown
        (populated via the \texttt{introspect} query);
  \item Choose columns to select from a checklist;
  \item Add filter conditions (column, operator, value);
  \item Add relation traversals (entity, foreign key, relation type);
  \item Set pagination and sort parameters.
\end{enumerate}
The form state is continuously serialised to a GraphQL query string
and displayed in a read-only textarea so users can see and copy the
equivalent raw query.  Submitting the form issues the query and
renders the JSON response in a collapsible tree view.

An \emph{examples} dropdown loads pre-built queries from the curated
catalog (Section~\ref{sec:catalog}), making it easy to explore public
datasets without any prior configuration.

% ===============================================================
\section{Database Catalog}
\label{sec:catalog}
% ===============================================================

Mozg ships a curated catalog of 63 pre-configured database and API
entries spanning several categories:

\begin{itemize}
  \item \textbf{Relational databases}: Chinook (SQLite music-store
        sample), Blog (SQLite), RNAcentral (PostgreSQL at EBI),
        RFAM (MySQL at EBI);
  \item \textbf{Graph databases}: Neo4j Movies demo;
  \item \textbf{REST / OpenAPI}: PokéAPI, Jikan (anime),
        Rick and Morty, SWAPI, GitHub, SpaceX, NASA,
        Open-Meteo, RestCountries, CoinGecko, MusicBrainz,
        OpenFoodFacts, Open Library, Art Institute of Chicago,
        Dog CEO, Fake Store, Open Brewery DB, Open Trivia, JSONPlaceholder;
  \item \textbf{Bioinformatics (NCBI family)}: NCBI E-utilities,
        GenBank, PubMed, RefSeq (NCBI Datasets v2), GEO;
  \item \textbf{Bioinformatics (protein \& structure)}: UniProt,
        PDB (Protein Data Bank), InterPro, BRENDA, HMDB,
        STRING protein interactions;
  \item \textbf{Bioinformatics (genome browsers \& pathways)}:
        Ensembl, Reactome, KEGG (via the dedicated \texttt{kegg} driver
        that converts KEGG's \texttt{text/plain} responses to JSON);
  \item \textbf{Bioinformatics (nucleotide archives)}: EMBL-EBI/ENA,
        DDBJ;
  \item \textbf{Bioinformatics (model organism databases)}: FlyBase
        (Drosophila), WormBase (C.~elegans), ZFIN (zebrafish);
  \item \textbf{BioCyc metabolic pathway databases (18 organisms)}:
        MetaCyc, EcoCyc, HumanCyc, YeastCyc, AraCyc, LeishCyc,
        TrypanocyC, plus eleven Tier~2 databases covering
        \textit{B.~subtilis}, \textit{M.~tuberculosis},
        \textit{P.~aeruginosa}, \textit{V.~cholerae},
        \textit{S.~enterica}, \textit{L.~monocytogenes},
        \textit{C.~crescentus}, \textit{C.~elegans},
        \textit{D.~melanogaster}, \textit{B.~burgdorferi}, and
        \textit{C.~jejuni}.
\end{itemize}

Each catalog entry defines the connection parameters, a list of known
entities (tables or API paths), the columns of interest, and the
cross-entity relations.  All public entries require no registration
and are suitable for immediate exploration.  Catalog data are loaded
lazily via CommonJS \texttt{require()} to keep startup time low.

% ===============================================================
\section{Evaluation}
\label{sec:evaluation}
% ===============================================================

\subsection{Functional Evaluation}

We exercised Mozg against the eight public data sources listed in
Table~\ref{tab:sources}.  For each source we verified:
(a)~schema introspection returns correct metadata,
(b)~a simple entity query returns expected rows,
(c)~at least one relation traversal produces a valid joined result.

All eight sources passed all three tests without modification to the
source data.  The cross-source join between RNAcentral (PostgreSQL)
and RFAM (MySQL) correctly joined RNA sequences with their family
annotations using a shared accession identifier.

\begin{table}[h]
\caption{Public data sources used in functional evaluation.}
\label{tab:sources}
\small
\begin{tabular}{llll}
\toprule
Source & Driver & Host & Data Domain \\
\midrule
RNAcentral & postgres & ebi.ac.uk & RNA sequences \\
RFAM & mysql & ebi.ac.uk & RNA families \\
Neo4j Movies & neo4j & demo.neo4jlabs.com & Film graph \\
PokéAPI & rest & pokeapi.co & Pokémon data \\
JSONPlaceholder & rest & jsonplaceholder.typicode.com & Fake posts \\
RestCountries & rest & restcountries.com & Country data \\
Open-Meteo & rest & api.open-meteo.com & Weather data \\
Chinook & sqlite3 & (local file) & Music store \\
\bottomrule
\end{tabular}
\end{table}

\subsection{Ontology Evaluation}

We ingested three ontologies of increasing complexity:
(1) the project's own \texttt{blog.ttl} (11 triples, 3 classes);
(2) the Friend of a Friend (FOAF) ontology~\cite{foaf}
(Turtle, $\sim$400 triples, 14 classes);
(3) a fragment of the Gene Ontology~\cite{go}
(OWL/XML, $\sim$50\,000 terms).

Parse times on a laptop with an Intel Core i7 were
$<$1\,ms, 12\,ms, and 3.1\,s respectively.  The SHA-256 cache
reduces repeat parse overhead to $<$0.1\,ms in all cases.  The
generated GraphQL SDL was syntactically valid in all three cases
as verified by graphql-js schema validation.

\subsection{Latency Analysis}

Table~\ref{tab:latency} reports round-trip latency (client-to-server
and back) for representative queries over the public sources.  All
measurements are medians over 20 requests from a residential broadband
connection (100 Mbit/s).  The dominant factor is network latency to
the remote data source, not Mozg's own processing overhead (which
accounts for $<$5\,ms in all cases measured locally).

\begin{table}[h]
\caption{Median round-trip latency (ms) for representative queries.}
\label{tab:latency}
\small
\begin{tabular}{lrr}
\toprule
Query & Latency (ms) & Rows \\
\midrule
Chinook Artists (local SQLite) & 8 & 275 \\
JSONPlaceholder /posts & 142 & 100 \\
PokéAPI /pokemon with types (relation) & 1\,840 & 20 \\
RNAcentral sequences (limit 50) & 290 & 50 \\
RFAM families (limit 20) & 340 & 20 \\
\bottomrule
\end{tabular}
\end{table}

The PokéAPI latency demonstrates the N+1 problem: fetching 20 Pokémon
plus their types requires 41 sequential HTTP requests (1 list + 20
detail + 20 type requests).  This is expected behaviour under the
current REST driver design and is documented as a deferred
performance issue.

% ===============================================================
\section{Discussion}
\label{sec:discussion}
% ===============================================================

\subsection{Design Trade-offs}

The ``slow by design'' philosophy accepts higher latency in exchange
for:
\begin{itemize}
  \item \textbf{Freshness}: every query reads live data with no
        staleness window;
  \item \textbf{Flexibility}: any change to source schemas is
        immediately reflected without pipeline rebuilds;
  \item \textbf{Simplicity}: no ETL pipelines, no data warehouse,
        no schema registry;
  \item \textbf{Minimal server state}: the server holds no persistent
        data beyond an in-memory connection cache and ontology cache,
        simplifying deployment.
\end{itemize}
These properties are valuable for exploratory analytics, scientific
workloads, and multi-tenant exploration tools where data ownership
prevents centralisation.

\subsection{Known Limitations and Deferred Issues}

Several security and performance issues are tracked but intentionally
deferred:

\begin{description}
  \item[SSRF] The ontology URL fetch in
    \texttt{src/ontology/index.js} accepts arbitrary URLs, enabling
    server-side request forgery.  Mitigation requires an allowlist and
    a per-request timeout.

  \item[Cypher injection] The Neo4j driver passes user-supplied
    identifiers (\texttt{from}, \texttt{orderBy}, \texttt{alias})
    directly into Cypher queries without escaping.

  \item[Path traversal] The static-file server does not decode
    percent-encoded path segments before resolving them against the
    document root.

  \item[Connection-cache collision] The in-memory connection cache is
    keyed on host and username without the password, so two users
    with the same username but different passwords on the same host
    could share a connection.

  \item[N+1 REST requests] The REST driver issues one HTTP request
    per row of the primary result to fetch related entities.
    Batching or parallelisation would significantly reduce latency.

  \item[BioCyc re-fetch] The BioCyc driver re-fetches the full
    class/instance list for every related entity ID rather than
    building an index.
\end{description}

\subsection{Scalability}

As a single-process Node.js application, Mozg's throughput is limited
by the event loop.  CPU-intensive ontology parsing blocks the loop;
large ontologies ($>$50\,000 triples) may cause noticeable latency
spikes.  Horizontal scaling is straightforward because the server
holds only in-memory caches that can be cold-started per instance;
queries are independently routable.

Query-level parallelism is partially exploited: within a single query,
sub-queries for independent relation branches could be issued
concurrently, though the current implementation issues them
sequentially.

\subsection{Future Work}

\begin{itemize}
  \item \textbf{Query optimisation}: pushdown filtering to source
        databases and parallel sub-query execution;
  \item \textbf{Schema registry}: persistent storage of discovered
        schemas and ontology results to accelerate repeated queries;
  \item \textbf{Authentication middleware}: JWT or OAuth2 bearer
        token validation to replace per-query credential passing;
  \item \textbf{Streaming results}: cursor-based pagination and
        incremental JSON streaming for large result sets;
  \item \textbf{SPARQL interface}: expose a SPARQL endpoint
        translating to the driver layer for triple-store
        compatibility;
  \item \textbf{WASM drivers}: compile lightweight database clients
        to WebAssembly to enable browser-side federation.
\end{itemize}

% ===============================================================
\section{Conclusion}
\label{sec:conclusion}
% ===============================================================

We presented Mozg, a ``slow by design'' federated query layer that
provides a single GraphQL interface over heterogeneous data sources.
Mozg's key contributions are:
\begin{enumerate}
  \item A \emph{driver model} that abstracts thirteen source types
        behind a uniform \texttt{executeQuery}/\texttt{introspect}
        interface, including a dedicated KEGG driver that converts
        \texttt{text/plain} responses to JSON;
  \item A \emph{relation-traversal engine} that joins entities across
        source boundaries at query time;
  \item An \emph{OWL ontology pipeline} that converts OWL TBox
        definitions into executable GraphQL schemas;
  \item A \emph{browser-based visual query builder} that democratises
        cross-database exploration for non-programmers;
  \item A \emph{curated catalog} of 63 publicly accessible data
        sources — including 19 major bioinformatics repositories
        and 18 BioCyc metabolic pathway databases — ready for
        immediate use.
\end{enumerate}

Mozg is designed for scenarios where data freshness, source
heterogeneity, and operational simplicity outweigh the need for
sub-millisecond latency.  By embracing this trade-off explicitly,
it enables a class of data-integration use cases that are poorly
served by traditional data warehouses or per-integration middleware.
The system is open-source, available at
\url{https://github.com/Skitionek/Mozg}, and can be launched
in a fully configured cloud environment via GitHub Codespaces in
under two minutes.

% ---------------------------------------------------------------
% Acknowledgements
% ---------------------------------------------------------------
\begin{acks}
The author thanks the open-source maintainers of GraphQL Yoga,
GraphQL Mesh, N3.js, Knex, and arangojs whose libraries form the
foundation of Mozg.

Portions of the code and documentation in this project were developed
with the assistance of \textbf{GitHub Copilot}, an AI-powered coding
assistant produced by GitHub and OpenAI\@.  The author reviewed all
AI-generated suggestions before inclusion and takes full responsibility
for the correctness and integrity of the work.
\end{acks}

% ---------------------------------------------------------------
% References
% ---------------------------------------------------------------
\bibliographystyle{ACM-Reference-Format}
\bibliography{mozg}

\end{document}
